\subsection{Terminal markings, phase shift and grounding}
\label{sec:na-terms}

The terminals of the higher-voltage winding are marked H1, H2, H3, those of
the lower-voltage winding X1, X2, X3, and the neutral X0 (H0 on a wye
primary)~[1]. In a voltage rating, ``$\times$'' joins two winding
sections that are connected in series or in parallel (240$\times$480\,V),
and ``/'' separates the line-to-line voltage from the line-to-neutral or
tap voltage (208Y/120\,V, 120/240\,V).

The line-to-neutral voltage of a wye winding is
\begin{equation}
  U_{ln} = \frac{U_{ll}}{\sqrt{3}}
  \label{eq:na-uln}
\end{equation}
which gives 277.1\,V for 480\,V, 138.6\,V for 240\,V and 120.1\,V for
208\,V.

Delta--delta and wye--wye connections have no phase shift. In delta--wye and
wye--delta connections the lower-voltage side lags the higher-voltage side
by 30\,$^{\circ}$~[1]; in the notation of IEC~60076-1 this is the
vector group Dyn1 or Yd1, Section~\ref{sec:de-vector}. Transformers can be
operated in parallel, or switched over without interruption, only when
their phase shifts are equal.

The secondary of a transformer with no direct electrical connection between
its windings is a separately derived system~[3, 250.30]. Its grounding is
decided at the transformer: the neutral or one phase is connected to the
equipment grounding conductor by the system bonding jumper and to a
grounding electrode, or the system is left ungrounded with ground
detection.
\FloatBarrier
